\begin{frame}{De l'automatisation des tâches à la menace de l'obsolescence des métiers techniques de la culture}
    Concept dont on discute depuis les années 1960, avec le développement de dispositifs techniques (en particulier l'informatique) \textrightarrow{} Automatisation de certaines tâches jusque-là assurées par des être humains.\\
    \vspace{1em}
Essayer de proposer de nouveaux services et donc une nouvelle manière de concevoir leur métier / le service ou l'objet qu'ils vendaient.\\
\textbf{Exemple:  }\\

\end{frame}

\begin{frame}{Exemple: l'informatisation du catalogue à la BnF}
    \textbf{Documentaire sur la BnF d'Alain Resnais, \textit{Toute la mémoire du monde}}\\
    \url{https://cours-edition-litterature-numerique-l2-92cac3.gitpages.huma-num.fr/img/Resnais_extrait.mp4}\\
    \vspace{1em}
    \pause
    \begin{itemize}
        \item Sous forme de fiches dans des catalogues imprimés;
        \item Pour Éric de Grolier (éditeur, chercheur, parfois considéré comme le fondateur des sciences de l'information en France), le documentaliste = « intermédiaire dont la fonction essentielle est de mettre en contact ceux qui ont besoin de savoir et ceux qui savent »
        \item Ailleurs on parle d’un « médiateur entre le document et l’utilisateur ».
        \item Gain d'autonomie énorme pour le lecteur
    \end{itemize}
    Reportage France 2: \url{https://cours-edition-litterature-numerique-l2-92cac3.gitpages.huma-num.fr/img/1765195014.mp4}
\end{frame}

\begin{frame}{Exemple: l'informatisation du catalogue à la BnF}
    \url{https://www.bnf.fr/fr/mediatheque/le-catalogue-collectif-de-france-quest-ce-que-cest}\\
    \vspace{1em}
    \textrightarrow{} Crainte de la dématérialisation (exemple Gallica).\\
    \textrightarrow{} On parle plutôt de dématérialisation.
\end{frame}